\clearpage\begingroup
\providecommand{\ObsSec}[2]{#1\mathord{\overset{\mathrm{s}}{,}}#2}
\providecommand{\ObsArc}[2]{#1\mathord{\overset{\prime\prime}{,}}#2}
\providecommand{\ObsJupiter}{\mathord{\tikz[baseline=-.1ex,x=1em,y=1em,line width=.045em]{
  \draw (.02,.48) .. controls (.04,.70) and (.33,.69) .. (.31,.46) .. controls (.29,.29) and (.14,.19) .. (.02,.13)--(.49,.13);
  \draw (.37,.64)--(.37,-.06);}}}
\providecommand{\ObsVesta}{\mathord{\tikz[baseline=-.1ex,x=1em,y=1em,line width=.045em]{
  \draw (.04,.02) rectangle (.40,.38);
  \draw (.22,.38)--(.22,.48) (.12,.63)--(.22,.48)--(.32,.63);}}}

\def\ObsJuno{\mathord{\tikz[baseline=-.2ex,x=1em,y=1em,line width=.045em]{
  \draw (.23,.72)--(.23,-.04);
  \draw (.05,.63)--(.41,.39) (.05,.39)--(.41,.63);
  \draw (.06,.13)--(.40,.13);}}}
\hypertarget{native-277-10}{}
\HistoricalHead{272}{L}{PRAKTISCHE UND SPHÄRISCHE ASTRONOMIE. NACHLASS UND BRIEFWECHSEL.}
\noindent [10. Seite] 134, 1811 Mai 11.

Heute Abend bei sehr heiterm Himmel wurde die Vesta aufgesucht: ein Stern 7. Gr[össe] $1^0$ östlich von dem Platze, der von Herrn Prof[essor] \textsc{Harding} berechneten Ephemeride wurde dafür angenommen und mit vier Sternen der \textit{Hist. Cél.} verglichen, die pag. 346 so vorkommen:
\begin{center}\begin{tabular}{c|rrr|rrr}
$a$&$16^{\mathrm h}$&$21^{\mathrm m}$&$\ObsSec{14}{21}$&$60^0$&$46'$&$38''$\\
$b$&16&22&$15\;60$&61&8&58\\
$c$&16&22&$35\;90$&60&54&48\\
$d$&16&26&$37\;89$&60&14&24
\end{tabular}\end{center}
\begin{center}\fontsize{9}{11}\selectfont\setlength{\tabcolsep}{2pt}
\begin{tabular}{rr|rr|rr|rr|rr}
\multicolumn{2}{c}{$a$}&\multicolumn{2}{c}{$b$}&\multicolumn{2}{c}{$c$}&\multicolumn{2}{c}{$\ObsVesta$}&\multicolumn{2}{c}{$d$}\\
$1^{\mathrm h}16^{\mathrm m}26^{\mathrm s}$&—&$16^{\mathrm m}49^{\mathrm s}$&$19^{\mathrm m}24^{\mathrm s}$&$17^{\mathrm m}16^{\mathrm s}$&$19^{\mathrm m}39^{\mathrm s}$&$21^{\mathrm m}22^{\mathrm s}$&$23^{\mathrm m}41^{\mathrm s}$&&\\
&&$28\quad32$&$29\quad28$&&&$32\quad21{,}5$&$34\quad27$&&\\
&&&&&&$49\quad25$&$50\quad22$&$50^{\mathrm m}25^{\mathrm s}$&$52^{\mathrm m}30^{\mathrm s}$
\end{tabular}\end{center}
\clearpage\SourceLeaf{278}\HistoricalHead{273}{R}{VERSCHIEDENE BEOBACHTUNGEN IN GÖTTINGEN.}
\noindent Hieraus beiläufig
\begin{center}\begin{tabular}{r lrrr}
$\ObsVesta\;10^{\mathrm h}13^{\mathrm m}$ M. Z.:&AR&$246^0$&$49'$&$21''$\\
&D&12&29&8.
\end{tabular}\end{center}
\medskip\noindent [11. Seite] 165, 1812 Febr. 8.

\{Von Mira Ceti hatte Herr Professor \textsc{Gauss} diesen Abend mit seinem Cometensucher gar keine Spur finden können\}\textsuperscript{[*)]}.
\medskip\noindent [12. Seite] 176, 1812 Juni 3.

Heute Abend bestätigte sich, dass der am 1. Juni für Juno gehaltene Stern bloss ein Fixstern gewesen sei, was schon am 2. Juni bemerkt war. Jedoch glaube ich die Juno heute Abend wirklich erkannt zu haben, da die Configuration der Sterne sich geändert hatte. Hienach nehme ich folgende geschätzte Positionen der Juno an:
\begin{center}\begin{tabular}{lr rrrr}
Jun.&1.&$283^0$&$17'$&$-4^0$&$54'$\\
&2.&283&8&$-4$&52\\
&3.&282&59&$-4$&50.
\end{tabular}\end{center}
\noindent\hspace*{\parindent}Die Ephemeride gibt also die Rectasc[ension] $14'$ zu klein und die Declination gut oder eine Kleinigkeit zu weit nördlich.

Was aus dem am 5. Mai gesehenen und nachher vermissten Sterne\textsuperscript{[**)]}, welcher $16'$ nördl[ich] und in einer ein Paar Min. grössern Rectascension stand, geworden sei, weiss ich nicht zu erklären.
\begingroup\renewcommand{\thefootnote}{[*)}\footnotetext{\textsc{Hardings} Handschrift.]}\endgroup
\begingroup\renewcommand{\thefootnote}{[**)}\footnotetext{Unter dem 5. Mai findet sich die Zeichnung einer] Karte zur Aufsuchung der Juno[, wo der später vermisste Stern durch grüne Farbe gekennzeichnet ist.]}\endgroup
\begin{center}\small
\rule{.81\textwidth}{.35pt}\\[.2em]
Monatliche Correspondenz, Band XXIII, S. 510—511, 1811 Mai.\\[-.2em]
\rule{.81\textwidth}{.35pt}
\end{center}
\begin{center}
[13.]\quad Auszug aus zwei Schreiben des Herrn Professor \textsc{Gauss}.\\[.4em]
Göttingen, am 25. Apr. u. 12. Mai 1811.
\end{center}
\noindent\hspace*{\parindent}\mbox{. . . . . .} Am 22. April war ich so glücklich, die Juno wirklich aufzufinden und ein Paarmal zu beobachten; am 23. fand ich nur eben Gelegenheit, die Richtigkeit meiner Vermuthung zu bestätigen; allein gestern sind wieder ein Paar Beobachtungen gelungen. Juno hat kaum die zehnte Grösse.
\HistoricalFoot{XI\textsubscript{1}.}{35}
\clearpage\SourceLeaf{279}\HistoricalHead{274}{L}{PRAKTISCHE UND SPHÄRISCHE ASTRONOMIE. BRIEFWECHSEL.}
\noindent Am 27. März hatte ich auch gewiss die Juno gesehen. Hier meine Beobachtungen:
\begin{center}\renewcommand{\arraystretch}{1.15}
\begin{tabular}{lr|rrr|rrr|rrr}
\hline\hline
\multicolumn{2}{c|}{1811}&\multicolumn{3}{c|}{Mittl. Z.}&\multicolumn{3}{c|}{AR. $\ObsJuno$}&\multicolumn{3}{c}{Südl. Decl.}\\\hline\hline
April&22&$9^{\mathrm h}$&$51^{\mathrm m}$&$35^{\mathrm s}$&$216^0$&$41'$&$50''$&$0^0$&$58'$&$16''$\\
&24&10&32&55&216&17&59&0&46&50
\end{tabular}\end{center}
\noindent\hspace*{\parindent}\textsc{Bodens} Ephemeride gibt also die gerade Aufsteigung jetzt um 7 Minuten zu gross, die Declinationen aber gut. Meine beobachteten Declinationen sind nicht ganz genau, da die Sterne nicht bequem standen. Heute Abend kommt Juno bis auf $1'$ einem Stern neunter Grösse nahe, der nicht in der \textit{Histoire céleste} steht und den ich Sie zu bestimmen bitte, da ich vielleicht heute Abend die Juno damit vergleiche. Seine scheinbare Stellung, so gut ich sie habe bestimmen können, ist:
\begin{center}
AR. $216^0\,6'\,59''$.\quad Decl. austr. $0^0\,38'\,39''$.
\end{center}
\noindent Nahe dabei steht ein Stern der \textit{Histoire céleste}, dessen Bestimmung ich ebenfalls wünschte
\begin{center}
AR. $216^0\,16'\,52''$.\quad Decl. austr. $0^0\,34'\,45''$.
\end{center}
\noindent\hspace*{\parindent}\mbox{. . . . . .} Seit meinem letzten Brief habe ich noch folgende Beobachtungen der Juno gemacht:
\begin{center}\renewcommand{\arraystretch}{1.15}
\begin{tabular}{lr|rrr|rrr|rrr}
\hline\hline
\multicolumn{2}{c|}{1811}&\multicolumn{3}{c|}{Mittl. Z.}&\multicolumn{3}{c|}{Ger. Aufst.}&\multicolumn{3}{c}{Decl. austr.}\\\hline\hline
April&25&$10^{\mathrm h}$&$19^{\mathrm m}$&$22^{\mathrm s}$&$216^0$&$6'$&$\ObsArc36$&$0^0$&$40'$&$\ObsArc36$\\
&26&10&17&30&215&54&$7{,}3$&0&33&$51{,}5$
\end{tabular}\end{center}
\noindent\hspace*{\parindent}Bisher hat es mir an Zeit gefehlt, Resultate daraus zu ziehen. Da die hiesigen Beobachtungen dieses Jahr vielleicht die einzigen sind, so würde es mir äusserst wichtig sein, die Position der verglichenen Sterne möglichst genau zu haben, besonders der zwei Ihnen schon das letztemal angezeigten, auf die sich die beiden letzten Declinationen gründen\textsuperscript{*)}.
\begingroup\renewcommand{\thefootnote}{*)}\footnotetext{Aus viertägigen am Passagen-Instrument und Quadranten gemachten Beobachtungen erhielt ich für den 19. Mai 1811 die scheinb. Position dieser Sterne:
\begin{center}\setlength{\tabcolsep}{3pt}
\begin{tabular}{llllll|lllll}
Nr.&I&AR.&$216^0$&$6'$&$\ObsArc{48}3$&Decl.&austr.&$0^0$&$38'$&$\ObsArc{13}0$\\
»&II&»&216&16&$48{,}5$&»&»&0&34&$21{,}6$.
\end{tabular}\end{center}
\noindent\hspace*{\parindent}Wir sind neugierig zu erfahren, ob noch irgend ein anderer Astronom die Juno beobachtet haben wird. Am hiesigen lichtstarken Passagen-Instrument war sie nicht ein einzigesmal sichtbar, auch wenn die Fäden ganz unerleuchtet blieben.\quad v. \textsc{L[indenau]}.}\endgroup
\clearpage\SourceLeaf{280}\HistoricalHead{275}{R}{VERSCHIEDENE BEOBACHTUNGEN IN GÖTTINGEN.}
\noindent\hspace*{\parindent}Den Cometen habe ich wegen Mondschein und schlechtem Wetter gar nicht suchen können. Gestern aber habe ich die Vesta beobachtet, aber nur flüchtig so eben reducirt:
\begin{center}
1811 Mai 11.\quad $10^{\mathrm h}13^{\mathrm m}$ M. Z.\quad $246^0\,49'$ AR.\quad $12^0\,29'$ süd. Ab.
\end{center}
\noindent\textsc{Hardings} Ephemeride weicht also $1^0\,18'$, die von \textsc{Bode} nur $22'$ in AR. ab. Der Planet ist übrigens leicht zu sehen, denn er hat 6.—7. Grösse.
\clearpage\endgroup
